%% Cover Letter — Physica A Submission
%% ELAPSE: Entropy-Linked Autonomous Propagation and Self-Erasure
%%
%% Compile: pdflatex ELAPSE_PhysicaA_CoverLetter.tex
%%
\documentclass[12pt,a4paper]{article}

\usepackage[margin=2.5cm]{geometry}
\usepackage{hyperref}
\usepackage{microtype}

\begin{document}

\noindent
\textbf{To the Editors of Physica A: Statistical Mechanics and its Applications}

\medskip

\noindent
Dear Editor-in-Chief,

\medskip

We submit for your consideration the manuscript:

\begin{center}
  \textbf{ELAPSE: Entropy-Linked Autonomous Propagation and Self-Erasure ---\\
  A Multi-Domain Ensemble Framework for Active Data Containment\\
  in Decentralised Networks}

  \smallskip
  Y.\ Sailaja, Chhavi Pareek, Hitarth Mehra, Harsh Patel\\
  RV College of Engineering, Bangalore, India
\end{center}

\bigskip

\noindent\textbf{Why this paper belongs in Physica A.}

This manuscript addresses a canonical problem in statistical mechanics: the
controlled containment of a diffusing payload on a complex network, viewed
as a phase-transition problem in the space of entropy threshold parameters.
The central result is a percolation-like phase transition at a critical
entropy threshold $H_c^*$, separating a subcritical regime (where
entropy-triggered deletion fires before data infects the bulk) from a
supercritical regime (where the network saturates before deletion can act).
The transition width is controlled by the algebraic connectivity $\lambda_2$
of the network Laplacian, directly connecting the result to Physica A's
long tradition of spectral graph theory and spreading phenomena.

\medskip

The paper contributes to three active areas in the Physica A community:

\begin{enumerate}
  \item \textbf{Information spreading and epidemic models on complex networks}
    (Pastor-Satorras \& Vespignani 2001; Boccaletti et al.\ 2006):
    We extend the standard SIR spreading framework to include active entropy-gated
    deletion, establishing scaling laws IEE~$\sim n^\alpha$ with
    topology-dependent exponents $\alpha \in [1.12, 1.40]$ across ER, BA,
    and WS topologies.

  \item \textbf{Percolation and phase transitions on random graphs}
    (Callaway et al.\ 2000; Newman et al.\ 2001):
    The critical entropy threshold $H_c^* = H_{\max} - C_0 e^{-2\alpha\lambda_2 T}$
    plays the role of a bond-percolation threshold, with the transition sharpness
    controlled by $\lambda_2$ (sharp for ER; diffuse for WS).

  \item \textbf{Stochastic diffusion processes on graphs}
    (Spectral graph theory; Fokker--Planck on networks):
    The ELAPSE framework is formalised as a stochastic differential equation
    with CIR-type noise, with well-posedness established via Yamada--Watanabe
    and a closed-form spectral IEE bound derived from the spectral gap.
\end{enumerate}

\medskip

\noindent\textbf{Summary of key findings.}

\begin{enumerate}
  \item A percolation-like entropy phase transition at $H_c^*$ with
    topology-dependent transition width $\sim 1/(2\alpha\lambda_2 T)$,
    empirically confirmed for ER, BA, and WS networks at $n \in \{100, 200, 500\}$.
  \item IEE scaling laws IEE~$\sim n^\alpha$ with $\alpha_{\rm ER} \approx 1.12$,
    $\alpha_{\rm BA} \approx 1.31$, $\alpha_{\rm WS} \approx 1.40$,
    consistent with the spectral prediction $\alpha \approx 1 + \beta$
    where $\lambda_2 \sim n^{-\beta}$.
  \item Five rigorous theoretical results including well-posedness
    (Yamada--Watanabe), a spectral IEE upper bound (via the Fiedler value),
    and exact convexity of the linearised IEE objective.
  \item A five-mechanism ensemble achieving $\geq 9\times$ IEE reduction
    over the uncontrolled baseline, validated on synthetic (ER/BA/WS)
    and real-world (Stanford SNAP Gnutella, Karate Club, Les Mis\'erables) networks.
  \item Connection to Jaynes' maximum entropy principle: the ensemble
    regulariser is the Gibbs distribution at a critical temperature
    identified by the Pareto elbow of the IEE--diversity trade-off.
\end{enumerate}

\medskip

\noindent\textbf{Manuscript details.}

\begin{itemize}
  \item Word count (main manuscript): approximately 11,500 words.
  \item Number of figures (main): 14 (Figs.\ 1--14), including 6 new
    simulation-driven figures (phase diagram, degree scaling, epidemic
    threshold, mechanism diversity, $n=1000$ scaling, extended real-world).
  \item Number of tables (main): 8.
  \item Number of references: 42.
  \item Supplementary material: 1 document (proofs, full tables).
  \item Highlights: 5 (submitted separately).
\end{itemize}

\medskip

\noindent\textbf{Suggested reviewers.}

We suggest the following experts with relevant expertise in the areas
addressed by this manuscript:

\begin{enumerate}
  \item \textbf{Prof.\ Romualdo Pastor-Satorras} (Universitat Polit\`ecnica
    de Catalunya) --- epidemic spreading on complex networks, SIR models,
    threshold phenomena.
  \item \textbf{Prof.\ Alessandro Vespignani} (Northeastern University) ---
    epidemic spreading, network dynamics, information diffusion.
  \item \textbf{Prof.\ Santo Fortunato} (Indiana University) --- community
    structure, network science, stochastic processes on graphs.
  \item \textbf{Prof.\ Sergey Dorogovtsev} (Universidade de Aveiro) ---
    critical phenomena on complex networks, percolation theory.
  \item \textbf{Prof.\ Marc Barth\'el\'emy} (CEA Saclay) --- dynamical
    processes on complex networks, spatial networks, spreading phenomena.
\end{enumerate}

\medskip

\noindent\textbf{Declarations.}

\begin{itemize}
  \item This manuscript is not currently under review elsewhere.
  \item The authors have no conflicts of interest to declare.
  \item All authors have reviewed and approved the manuscript.
  \item Data and code availability: simulation code and data will be made
    available on acceptance.
\end{itemize}

\medskip

We believe this paper makes a substantive contribution to the statistical
physics of controlled diffusion on complex networks, and that Physica A is
the most appropriate venue for this work.  We look forward to the review
process.

\bigskip

\noindent Sincerely,

\bigskip\bigskip

\noindent Dr.\ Y.\ Sailaja (Corresponding Author)\\
Department of Mathematics\\
RV College of Engineering\\
Bangalore 560059, India\\
Email: sailaja@rvce.edu.in

\end{document}
